\DocumentMetadata{
  pdfversion=2.0,
  pdfstandard=ua-2,
  lang=en-US,
  tagging=on,
  tagging-setup={math/setup=mathml-SE,tabsorder=structure}
}
\documentclass[11pt,letterpaper]{article}
\usepackage[margin=0.68in,includefoot]{geometry}
\usepackage{newcomputermodern}
\usepackage{amsmath,amscd,mathtools}
\usepackage{tikz,adjustbox}
\providecommand{\square}{\mdlgwhtsquare}
\usepackage[hidelinks]{hyperref}
\hypersetup{
  pdftitle={Complex Analysis PhD Exam, May 2006},
  pdfauthor={University of Florida Department of Mathematics},
  pdfsubject={Graduate mathematics examination},
  pdfdisplaydoctitle=true,
  bookmarksopen=true
}
\setcounter{secnumdepth}{0}
\setlength{\parindent}{0pt}
\setlength{\parskip}{0.28em}
\setlength{\emergencystretch}{3em}
\setlength{\leftmargini}{2.15em}
\setlength{\leftmarginii}{2.65em}
\setlength{\leftmarginiii}{3em}
\pagestyle{empty}
\raggedbottom
\allowdisplaybreaks
\NewTaggingSocketPlug{block/list/label}{examproblem}{%
  \tagstructbegin{tag=Lbl}\tagstructend%
  \tagstructbegin{tag=\LBody}%
  \tagstructbegin{tag=\ProblemHeadingTag,title-o={\ProblemHeadingText}}%
  \tagmcbegin{tag=\ProblemHeadingTag,actualtext={\ProblemHeadingText}}%
  #2%
  \tagmcend%
  \tagstructend%
}
\newenvironment{examproblems}[2]{%
  \def\ProblemHeadingTag{#2}%
  \AssignTaggingSocketPlug{block/list/label}{examproblem}%
  \begin{enumerate}%
  \setcounter{enumi}{#1}%
  \renewcommand{\labelenumi}{\textbf{\theenumi.}}%
  \setlength{\topsep}{0.35em}%
  \setlength{\partopsep}{0pt}%
  \setlength{\itemsep}{0.48em}%
  \setlength{\parsep}{0.18em}%
}{\end{enumerate}}
\newenvironment{examsubparts}{%
  \AssignTaggingSocketPlug{block/list/label}{default}%
  \begin{enumerate}%
  \setlength{\topsep}{0.18em}%
  \setlength{\partopsep}{0pt}%
  \setlength{\itemsep}{0.16em}%
  \setlength{\parsep}{0.1em}%
}{\end{enumerate}}
\newenvironment{examinstructions}{%
  \par\begingroup\itshape
}{%
  \par\endgroup\vspace{0.18em}
}
\makeatletter
\renewcommand\subsection{\@startsection{subsection}{2}{0pt}%
  {-1.25ex plus -.25ex minus -.1ex}%
  {0.55ex plus .1ex}%
  {\normalfont\large\bfseries}}
\renewcommand{\@maketitle}{%
  \newpage\null\vskip -1em%
  {\footnotesize\bfseries
    \href{https://www.ufl.edu/}{University of Florida}, \href{https://math.ufl.edu/}{Department of Mathematics}\par}%
  \vskip -0.5em%
  \tagstructbegin{tag=H1,title={Complex Analysis PhD Exam, May 2006}}%
  \tagpdfparaOff%
  \tagmcbegin{tag=H1}%
    {\LARGE\bfseries\href{https://gma.math.ufl.edu/exams/complex-analysis-phd/}{Complex Analysis PhD Exam}, May 2006\par}%
  \tagmcend%
  \tagpdfparaOn%
  \tagstructend%
  \vskip 0.25em%
  \tagmcbegin{artifact}%
    \hrule height 1pt%
  \tagmcend%
  \vskip 1em%
}
\makeatother
\title{Complex Analysis PhD Exam, May 2006}
\author{}
\date{}
\begin{document}
\maketitle
\thispagestyle{empty}
\begin{examinstructions}
Give complete proofs and computations. Partial credit will be given where justified.
\end{examinstructions}
\begin{examproblems}{0}{H2}
\def\ProblemHeadingText{Problem 1.}
\item
Let \(0<a<1\). Evaluate the integral
\[
\int_{-\infty}^{\infty}\frac{e^{ax}}{1+e^x}\,dx,
\]
 where \(a\in(0,1)\).
\def\ProblemHeadingText{Problem 2.}
\item
This problem has two parts.
\begin{examsubparts}
\item[\textnormal{(a)}]
Prove that
\[
\cos \pi z=\prod_{n=1}^{\infty}\left(1-\frac{4z^2}{(2n-1)^2}\right).
\]
\item[\textnormal{(b)}]
Deduce from (a) that
\[
\tan \pi z=\frac{8}{\pi}\sum_{n=1}^{\infty}\frac{z}{(2n-1)^2-4z^2}.
\]
\end{examsubparts}
\def\ProblemHeadingText{Problem 3.}
\item
Let~\(G\) be a region and \(u:G\to\mathbb{R}\) be harmonic. Let~\(v\) be a harmonic conjugate to~\(u\) in~\(G\). Show that the product \(u\cdot v\) is harmonic.
\def\ProblemHeadingText{Problem 4.}
\item
This problem has two parts.
\begin{examsubparts}
\item[\textnormal{(a)}]
Show that if~\(f\) is analytic on the closed disk \(\overline{B(0;1)}\), then
\[
|f(a)|^2\leq \frac{1}{\pi R^2}\int_0^{2\pi}\int_0^R|f(a+re^{i\theta})|^2r\,dr\,d\theta.
\]
\item[\textnormal{(b)}]
Let~\(G\) be a region and \(M\in\mathbb{R}\) be fixed. Let
\[
\mathcal{F}=\left\{f\in H(G)\mathrel{\bigg|}\int_G|f(z)|^2\,dx\,dy\leq M\right\}.
\]
 Show that~\(\mathcal{F}\) is a normal family.
\end{examsubparts}
\def\ProblemHeadingText{Problem 5.}
\item
Let \(f(z)=\sum_{n=0}^{\infty}a_nz^n\) and \(g(z)=\sum_{n=0}^{\infty}b_nz^n\) be analytic on \(\overline{B(0;1)}\). Let
\[
F(z)=\frac{1}{2\pi i}\int_{\partial B(0;1)}\frac{f(w)}{w}g\left(\frac{z}{w}\right)\,dw,\qquad z\in B(0;1),
\]
 with the contour taken with counterclockwise orientation. Show that
\[
F(z)=\sum_{n=0}^{\infty}a_nb_nz^n,\qquad z\in B(0;1).
\]
\def\ProblemHeadingText{Problem 6.}
\item
Let \(G=\{z\in\mathbb{C}\mid |z|<1\text{ and }|2z-1|>1\}\), and let \(f\in H(G)\).
\begin{examsubparts}
\item[\textnormal{(a)}]
Must there exist a sequence of polynomials \(P_n\) such that \(P_n\to f\) uniformly on compact subsets of~\(G\)? Explain your reasoning.
\item[\textnormal{(b)}]
Must there exist such a sequence which converges to~\(f\) uniformly on~\(G\)? Explain your reasoning.
\item[\textnormal{(c)}]
Must there exist such a sequence which converges to~\(f\) uniformly on~\(G\), if~\(f\) is now required to be analytic in some neighborhood of \(\overline{G}\)? Explain your reasoning.
\end{examsubparts}
\def\ProblemHeadingText{Problem 7.}
\item
Let~\(G\) be a region and~\(A\) be a discrete and relatively closed subset of~\(G\). Assume that \(f\in H(G\setminus A)\) is injective.
\begin{examsubparts}
\item[\textnormal{(a)}]
Prove that no point of~\(A\) can be an essential singularity of~\(f\).
\item[\textnormal{(b)}]
Prove that if \(a\in A\) is a pole for~\(f\), then~\(a\) is a pole of order 1.
\end{examsubparts}
\def\ProblemHeadingText{Problem 8.}
\item
Let~\(f\) be analytic on \(\overline{B(0;1)}\). If \(|f(z)|<1\) on \(\partial B(0;1)\), prove that there exists a unique \(z_0\in B(0;1)\) (counting multiplicity) such that \(f(z_0)=z_0\). If \(|f(z)|\leq 1\) on \(\partial B(0;1)\), does the same conclusion hold? Justify your assertion.
\end{examproblems}
\end{document}
